\chapter{The Plane Cubic and Its Group Law}
\label{ch:plane-cubic-group}

Let \(E\) be the smooth cubic from \cref{ch:weierstrass-embedding}, with identity
point \(O=[0:1:0]\).  A line in \(\CP^2\) meets a cubic in three points counted
with multiplicity.  This is the only geometric input needed for the chord--
tangent rule.

\begin{definition}
Given \(P,Q\in E\), draw the line through \(P\) and \(Q\), using the tangent line
if \(P=Q\).  Let \(R\) be the third intersection point with \(E\).  Define
\[
  P\boxplus Q=\ominus R,
\]
where \(\ominus(x,y)=(x,-y)\) in the affine chart and \(\ominus O=O\).
\end{definition}

\begin{center}
\resizebox{0.62\linewidth}{!}{\classicChordDiagram}
\end{center}

The reflection step is often the mysterious part.  It is not arbitrary: a line
through \(P,Q,R\) cuts a principal divisor
\[
  [P]+[Q]+[R]-3[O].
\]
By Abel's theorem, this means \(P+Q+R=O\) in the group transported from the torus.
Therefore the group sum of \(P\) and \(Q\) is \(-R\).

\begin{proposition}
\label{prop:line-divisor}
If a projective line \(L\) meets \(E\) in \(P,Q,R\), counted with multiplicity,
then \([P]+[Q]+[R]-3[O]\) is principal.
\end{proposition}

\begin{proof}[Proof sketch]
The quotient of the linear form defining \(L\) by the linear form defining the
line at infinity is a meromorphic function on \(E\).  Its zeros are the
intersections with \(L\), and its poles are the triple intersection with the line
at infinity at \(O\).
\end{proof}

\section*{Exercises}
\addcontentsline{toc}{section}{Exercises}

\exerciseintent{See tangent lines as limiting secants in the nonvertical affine case.}
\begin{exercise}
Let \(P\) and \(Q\) be nearby affine points on a smooth cubic, and assume the
tangent at \(P\) is not vertical.  Let \(Q\to P\).  Compute the limiting slope
of the chord line in local affine coordinates.  State why the limiting line has
intersection multiplicity \(2\) at \(P\).
\end{exercise}

\exerciseintent{Derive the chord formula from substitution and Vieta's formula.}
\begin{exercise}
For \(y^2=4x^3-g_2x-g_3\), derive the formula
\[
  x_{P\boxplus Q}=\frac{m^2}{4}-x_P-x_Q
\]
when the chord has slope \(m\).
Do this by substituting \(y=mx+b\), comparing the \(x^2\)-coefficient of the
resulting cubic, and applying Vieta's formula.
\end{exercise}

\exerciseintent{Compute the vertical-line case that gives inverses on the cubic.}
\begin{exercise}
If \(P=(x,y)\), compute the intersections of the vertical line \(X=xZ\) with
the cubic.  Show that they are \(P\), \(\ominus P=(x,-y)\), and \(O\), counted
with multiplicity when necessary.
\end{exercise}

\exerciseintent{Compute the line divisor directly from the meromorphic function \(\ell/Z\).}
\begin{exercise}
Let \(\ell(X,Y,Z)=0\) be a projective line meeting \(E\) in \(P,Q,R\), counted
with multiplicity, and let \(h=Z\) define the line at infinity.  Compute the
divisor of the meromorphic function \(\ell/h\) on \(E\): identify its zeros from
\(\ell|_E=0\), identify its poles from \(Z|_E=0\), and use the fact that the
line at infinity meets the Weierstrass cubic triply at \(O=[0:1:0]\).  Conclude
that
\[
  \divisor(\ell/h)=[P]+[Q]+[R]-3[O].
\]
\end{exercise}

\exerciseintent{Turn the line divisor into the group equation \(P+Q+R=O\) by Abel sums.}
\begin{exercise}
Assume Abel's theorem and the transported group law on \(E\).  Use the line
divisor \([P]+[Q]+[R]-3[O]\) to compute the Abel sum and prove that
\(P+Q+R=O\).
\end{exercise}
